\documentclass[11pt,a4paper]{article}
\usepackage[utf8]{inputenc}
\usepackage[T1]{fontenc}
\IfFileExists{lmodern.sty}{\usepackage{lmodern}}{}
\IfFileExists{french.ldf}{\usepackage[french]{babel}}{\usepackage[english]{babel}}
\usepackage{amsmath,amssymb,amsfonts,amsthm}
\usepackage{geometry}
\geometry{top=2cm, bottom=2cm, left=2.2cm, right=2.2cm}
\usepackage{booktabs,array,multirow}
\usepackage{enumitem}
\usepackage{xcolor}
\usepackage{tcolorbox}
\tcbuselibrary{skins,breakable}
\usepackage{tikz}
\usetikzlibrary{arrows.meta,positioning,shapes.geometric,calc}
\usepackage[colorlinks=true,linkcolor=blue!60!black]{hyperref}
\newtcolorbox{propbox}[1]{colback=blue!5!white,colframe=blue!75!black,fonttitle=\bfseries,title={#1},breakable,boxrule=0.8pt}
\newtcolorbox{rembox}[1]{colback=orange!5!white,colframe=orange!80!black,fonttitle=\bfseries,title={#1},breakable,boxrule=0.8pt}
\newtcolorbox{verifbox}{colback=green!4!white,colframe=green!45!black,fonttitle=\bfseries,title=Vérification algorithmique (Python),breakable,boxrule=0.6pt}
\newcommand{\Xt}{\tilde{X}}
\newcommand{\Yt}{\tilde{Y}}
\setlength{\parindent}{0pt}
\setlength{\parskip}{4pt}
\title{\textbf{Université Akli Mohand Oulhadj -- Bouira}\\
\large Master Recherche Opérationnelle (RO) -- 2025/2026\\
\Large \textbf{Théorie des Jeux : Corrigé commenté du Test (QCM, 1h)}}
\author{\textbf{Référence de cours :} Chapitre 1 (jeux matriciels, PL, dualité)\\ Chapitre 2 (jeux antagonistes, Nash)}
\date{}

\newcommand{\bon}[1]{\textcolor{green!50!black}{\textbf{#1}}}
\newcommand{\faux}[1]{\textcolor{red!70!black}{#1}}

\begin{document}
\maketitle
\hrule
\vspace{0.3cm}

Chaque question peut avoir de 0 à 4 bonnes réponses. Les réponses justes sont en \bon{vert gras}. Lorsqu'une proposition est discutable, c'est signalé par le symbole $(\ast)$ avec l'argument.

\subsection*{Grille de réponses}
\begin{center}
\renewcommand{\arraystretch}{1.25}
\begin{tabular}{|c|c||c|c||c|c||c|c|}
\hline
Q & Réponses & Q & Réponses & Q & Réponses & Q & Réponses\\\hline
1 & a, c, d & 6 & a, d $(\ast)$ & 11 & b, c & 16 & b\\
2 & d & 7 & a, c, d & 12 & a, b & 17 & c\\
3 & a, b & 8 & a, b & 13 & a & 18 & a\\
4 & c, d & 9 & b & 14 & b, d & 19 & c, d\\
5 & c, d & 10 & a, d & 15 & a, c & 20 & b, c $(\ast)$\\\hline
\end{tabular}
\end{center}

\section*{Bloc 1 : Notions générales (Q1--Q8)}

\paragraph{Q1. La théorie des jeux est un ensemble de méthodes mathématiques\dots}
\bon{(a)} où interviennent des règles de décision ; \faux{(b)} « tactiques » n'est pas un concept de la théorie ; \bon{(c)} on manipule des notions abstraites de \emph{stratégies} ; \bon{(d)} elle résout des problèmes concurrentiels (« situations de conflit », Chap.~1 p.~1).

\paragraph{Q2. Dans un duel, le gain ou la perte d'un joueur dépend\dots}
\faux{(a)}, \faux{(b)} : faux, le résultat dépend des choix \emph{des deux} joueurs (une issue $=$ résultat des choix de tous les joueurs). \faux{(c)} : un duel est à somme \emph{nulle}. \bon{(d)} : de l'adresse et de l'intelligence (la qualité des décisions) des joueurs.

\paragraph{Q3. Le jeu du duel est représenté par\dots}
\bon{(a)} le gain d'un joueur et \bon{(b)} la perte d'un joueur : dans un jeu à deux personnes à somme nulle, $f_1+f_2=0$, « on mentionne juste le gain du joueur 1 » (Chap.~1 p.~3), qui est aussi la perte du joueur 2. \faux{(c)}, \faux{(d)} : non (un duel a deux joueurs).

\paragraph{Q4. La forme normale correspond\dots}
\faux{(a)}, \faux{(b)} : non. \bon{(c)} un tableau à double entrée (lignes $\times$ colonnes) ; \bon{(d)} un tableau des gains (ou pertes) d'un joueur : la matrice $A=(a_{ij})$ du jeu matriciel.

\paragraph{Q5. Matrice $m\times n$ des joueurs A et B.}
Le joueur 1 (A) choisit les \emph{lignes} : $X=\{1,\dots,m\}$ ; le joueur 2 (B) les \emph{colonnes} : $Y=\{1,\dots,n\}$ (Chap.~1 p.~4). \bon{(c)}, \bon{(d)}.

\paragraph{Q6. Le jeu se présente comme suit\dots}
A choisit une \textbf{ligne}, B une \textbf{colonne} : \faux{(b)}, \faux{(c)} sont exclues. \bon{(a)} décrit le raisonnement \emph{maximin} de A (pour chaque ligne, B répond par la colonne la pire pour A). $(\ast)$ \bon{(d)} décrit le raisonnement \emph{minimax} de B (pour chaque colonne, A répond par la meilleure ligne). Les deux attribuent correctement lignes/colonnes. \emph{Attention :} dans la forme normale, les choix sont simultanés ; si le correcteur n'attend qu'une réponse, c'est (a).

\paragraph{Q7. Valeur du jeu (somme nulle).}
\bon{(a)} c'est $a_{i^*j^*}$, à l'intersection de la ligne choisie par A et de la colonne choisie par B (point selle). \faux{(b)} inverse les rôles. \bon{(c)} si $v>0$, c'est ce que B paie à A ; \bon{(d)} si $v<0$, A paie $|v|$ à B (« peut être »).

\paragraph{Q8. Dans un jeu à stratégie simple (pure), les joueurs doivent\dots}
\bon{(a)} être prudents : le critère maximin/minimax est un critère de prudence (gain \emph{garanti}) ; \bon{(b)} être intelligents (rationnels, ils connaissent $X,Y,f$). \faux{(c)} prendre des risques : non. \faux{(d)} coopérer : non, le jeu est non coopératif et antagoniste.

\section*{Bloc 2 : Stratégies pures et point selle (Q9--Q13)}

\begin{propbox}{Rappel}
$\alpha=\max_i\min_j a_{ij}$ (maximin, joueur 1 : min de chaque ligne, puis max) ; $\beta=\min_j\max_i a_{ij}$ (minimax, joueur 2 : max de chaque colonne, puis min). Toujours $\alpha\le\beta$. Point selle $\iff\alpha=\beta=a_{i^*j^*}$, élément \emph{minimal dans sa ligne et maximal dans sa colonne} ; sa valeur est la valeur du jeu.
\end{propbox}

\paragraph{Q9.} \bon{(b)} le premier joueur choisit le maximin et le second le minimax. \faux{(a)} trop vague ; \faux{(c)}, \faux{(d)} inversent les rôles.

\paragraph{Q10. Joueur A.} \bon{(a)} il repère le plus petit $a_{ij}$ de chaque ligne puis choisit la ligne où cette valeur est la plus grande ; \bon{(d)} c'est le maximin. \faux{(b)}, \faux{(c)} décrivent le joueur B.

\paragraph{Q11. Joueur B.} \bon{(c)} il repère le plus grand $a_{ij}$ de chaque colonne puis choisit la colonne où cette valeur est la plus petite ; \bon{(b)} c'est le minimax. \faux{(a)}, \faux{(d)} décrivent A.

\paragraph{Q12.} \bon{(a)} la valeur du point selle est la valeur du jeu ($v=a_{i^*j^*}$) ; \bon{(b)} c'est aussi une valeur d'équilibre (« on obtient un équilibre quand $\alpha=\beta$ », Chap.~1 p.~7). \faux{(c)}, \faux{(d)} : non.

\paragraph{Q13.} \bon{(a)} un élément à la fois \emph{le plus petit de sa ligne et le plus grand de sa colonne} (propriété 1, p.~8). \faux{(b)}, \faux{(c)}, \faux{(d)} : non. Exemple du cours : $a_{12}=a_{14}=3$ dans $\begin{pmatrix}5&3&4&3\\7&2&0&-2\\0&-1&-4&2\end{pmatrix}$.

\section*{Bloc 3 : Programmation linéaire et dualité (Q14--Q20)}

\begin{propbox}{Rappel : correspondance primal / dual (forme canonique)}
\centering
\begin{tabular}{l|l}
Primal & Dual\\\hline
$\max c^Tx$ & $\min b^Tu$\\
$Ax\le b$ ($m$ contraintes) & $A^Tu\ge c$ ($n$ contraintes)\\
$x\ge0$ ($n$ variables) & $u\ge0$ ($m$ variables)
\end{tabular}

\raggedright
Mêmes données ($A,b,c$), réarrangées : la matrice est \textbf{transposée}, seconds membres et coûts échangés, max $\leftrightarrow$ min, $\le\ \leftrightarrow\ \ge$. Pour un jeu : $(III)$ $\min\sum p_i$, $\sum_ia_{ij}p_i\ge1$ et $(IV)$ $\max\sum q_j$, $\sum_ja_{ij}q_j\le1$ sont duaux (Chap.~1 p.~14--15).
\end{propbox}

\paragraph{Q14. Différences entre primal et dual.}
\faux{(a)} les données chiffrées sont les mêmes (seulement réarrangées). \bon{(b)} la forme d'écriture diffère ; \faux{(c)} « et les données » : faux ; \bon{(d)} l'objectif diffère (max $\leftrightarrow$ min).

\paragraph{Q15. Primal à maximiser.}
\bon{(a)} le dual est à minimiser. \faux{(b)} en forme canonique, les contraintes d'un max sont $\le$. \bon{(c)} les contraintes sont toutes de signe $\le$, sauf les contraintes de positivité $x\ge0$ (le symbole « $\le$ » manque dans l'énoncé). \faux{(d)} les variables sont $\ge0$.

\paragraph{Q16.} Primal avec contraintes $\ge$ (un min canonique) $\Rightarrow$ dual avec contraintes \bon{(b)} $\le$. \faux{(a)}, \faux{(c)}, \faux{(d)} : non (en PL on n'utilise pas d'inégalités strictes).

\paragraph{Q17.} Les matrices des coefficients sont \bon{(c)} transposées l'une de l'autre. \faux{(a)}, \faux{(b)}, \faux{(d)} : non.

\paragraph{Q18.} Le théorème des écarts complémentaires est un corollaire \bon{(a)} du théorème fondamental de la dualité (égalité des valeurs optimales primal/dual). \faux{(b)}, \faux{(c)}, \faux{(d)} : non.

\paragraph{Q19. Avantages du T.E.C.}
\bon{(c)} obtenir la solution du dual à partir de celle du primal, et \bon{(d)} inversement. C'est exactement l'usage du cours (p.~29) : de $p_1,p_2>0$ on déduit des contraintes duales saturées, et de $4p_2>1$ on déduit $q_3=0$. \faux{(a)}, \faux{(b)} : ce ne sont pas ses avantages principaux.

\paragraph{Q20. Objectifs de la dualité.}
\bon{(b)} simplification des calculs (résoudre le programme le plus simple ; ex. un primal à 2 variables résolu graphiquement) ; \bon{(c)} interprétation économique (prix fictifs/marginaux). \faux{(a)} non. $(\ast)$ \faux{(d)} le développement mathématique de la PL est un apport théorique, pas un objectif principal ; certains correcteurs l'acceptent.

\begin{rembox}{Conseils pour le QCM}
\begin{itemize}[leftmargin=*]
\item Repérer les \emph{inversions de rôles} (ligne/colonne, maximin/minimax, max/min) : c'est le piège le plus fréquent (Q6, Q9--Q11).
\item « Uniquement », « seulement », « toutes » rendent souvent une proposition fausse (Q2).
\item En dualité, retenir : \emph{transposée, max $\leftrightarrow$ min, $\le\ \leftrightarrow\ \ge$, mêmes données}.
\end{itemize}
\end{rembox}

\begin{verifbox}
Les items calculatoires sous-jacents (définition du point selle, $\alpha$, $\beta$, exemple du cours p.~8 avec points selle $(1,2)$ et $(1,4)$, $v=3$ ; dualité des programmes $(III)$/$(IV)$ et écarts complémentaires sur les exemples des TD) sont recontrôlés dans \texttt{verification\_TDJ.py}. Les items purement conceptuels (Q1--Q4, Q8, Q20) relèvent de l'interprétation du cours et sont justifiés ci-dessus.
\end{verifbox}

\end{document}
